\chapter{Operative Vaginal Delivery}

\section*{Overview}
Operative vaginal delivery uses forceps or a vacuum extractor to assist delivery of a fetus through the vagina when spontaneous delivery cannot be achieved safely. The exact approach, setting, and preparation depend on the indication, anatomy, and the patient's health.

\section*{Alternative Names}
Forceps delivery; vacuum-assisted delivery; instrumental delivery; ventouse delivery

\section*{Principle}
Forceps cradle the fetal head with two curved blades; a vacuum extractor applies suction to the scalp. Traction during contractions assists descent of the head, shortening the second stage when indicated.

\section*{History}
Obstetric forceps were developed in the 17th century and held as family secrets by the Chamberlen family. The vacuum extractor was developed in the mid-20th century as an alternative.

\section*{Why It Is Done}
Performed for prolonged second stage, fetal distress, maternal exhaustion, or the need to shorten pushing for medical conditions such as cardiac or hypertensive disease.

\section*{Is This a Primary or Secondary Test or Procedure}
Secondary obstetric intervention when spontaneous delivery is not progressing safely.

\section*{How It Is Performed}
Performed after full cervical dilation with ruptured membranes and an engaged head. Regional anesthesia is preferred. Forceps blades are applied to the fetal head and traction is applied with contractions; a vacuum cup is placed on the occiput and gentle traction applied. Delivery is then completed.

\section*{Preparation}
Assessment of fetal position, station, and presentation; bladder emptying; consent; and readiness for emergency cesarean if delivery fails.

\section*{Contraindications and Precautions}
Non-engaged head, unknown fetal position, face presentation, fetal macrosomia with a small pelvis, and fetal bone disorders are contraindications. Precaution: failed instrumentation mandates cesarean delivery.

\section*{Risks}
Maternal perineal or vaginal trauma, third- or fourth-degree lacerations, hemorrhage, and fetal scalp injury, cephalohematoma, or rarely intracranial hemorrhage.

\section*{Aftercare and Recovery}
The perineum is examined and repaired, and the mother and baby are monitored. Pain and perineal care follow standard postpartum care.

\section*{Outcome and Follow-up}
Successful vaginal delivery with minimal trauma confirms a good outcome. Neonatal scalp findings typically resolve.

\section*{Expected Results}
A safe delivery of a healthy infant with controlled maternal trauma and a favorable neonatal course.

\section*{Follow-up}
Postpartum perineal care, newborn examination, and follow-up for any maternal or neonatal complications.

\section*{When to Seek Medical Care}
Follow the procedure team's specific instructions. Seek urgent medical assessment for severe or worsening pain, uncontrolled bleeding, fever or signs of infection, trouble breathing, chest pain, fainting, new weakness or confusion, inability to urinate, or any rapidly worsening symptom. The appropriate warning signs depend on the procedure and the patient's medical condition.

\section*{References}
\begin{enumerate}
  \item MedlinePlus. Forceps delivery. U.S. National Library of Medicine; 2025.
  \item Current Medical Diagnosis and Treatment. McGraw-Hill; 2025.
\end{enumerate}

\clearpage
